% =====================================================================
%  PLANTILLA: memoria de práctica / trabajo (portada completa + índice)
%  Máster Universitario en Ingeniería Informática (IA) - UPO 2026/2027
%
%  Compilar: pdflatex -> bibtex -> pdflatex -> pdflatex
%            (o simplemente: latexmk -pdf memoria.tex)
%  Opciones de la clase: portada | cabecera | sinindice | twoside
% =====================================================================
\documentclass{upo-master}

% ---------- Datos del documento ----------
\asignatura{Nombre de la asignatura}
\tipodocumento{Práctica 1}            % Práctica N, Actividad N, Memoria, Trabajo final...
\title{Título de la práctica}
\subtitulo{Subtítulo opcional}         % déjalo vacío o bórralo si no hace falta
\author{Javier Gómez López}
% \correo{usuario@alu.upo.es}
% \profesor{Nombre del profesor}
% \grupo{Grupo 1}
% \tituloCorto{P1 - Título corto}       % texto de la cabecera derecha
% \date{15 de diciembre de 2026}         % por defecto, la fecha de compilación
% \rutaLogo{../../Plantilla LaTeX/img/logo-upo}  % si el logo no está en img/

% ---------- Uso de IA (OBLIGATORIO revisarlo en cada entrega) ----------
% Una línea por herramienta. Si no declaras ninguna, el documento
% indicará que NO se ha usado IA generativa.
\herramientaIA{Claude (Anthropic)}{Claude Opus 5.5}{Apoyo en la redacción y revisión de las explicaciones, estructura de la memoria y código \LaTeX.}
% \herramientaIA{GitHub Copilot}{---}{Autocompletado de código.}
\usoIA{Describe aquí con más detalle cómo se ha usado la IA: qué partes,
qué se ha revisado y qué se ha hecho sin ella.}

\begin{document}

\maketitle

% ======================= INTRODUCCIÓN =======================
\section{Introducción y objetivos}

Breve contexto del problema y de lo que se pide en el enunciado.

\begin{enunciado}
Resumen o cita textual del enunciado de la práctica.
\end{enunciado}

Los objetivos son:
\begin{enumerate}
  \item Primer objetivo.
  \item Segundo objetivo.
\end{enumerate}

% ======================= FUNDAMENTOS =======================
\section{Fundamentos teóricos}

\begin{concepto}{Concepto clave}
Definición breve del concepto principal que se usa en el trabajo.
\end{concepto}

\begin{definicion}
Una definición formal numerada. Por ejemplo, el riesgo empírico de un
modelo $f_{\vect{\theta}}$ sobre $n$ ejemplos es
\begin{equation}
  \hat{R}(\vect{\theta}) = \frac{1}{n}\sum_{i=1}^{n} \ell\bigl(f_{\vect{\theta}}(\vect{x}_i), y_i\bigr),
  \qquad
  \vect{\theta}^{*} = \argmin_{\vect{\theta}\in\R^d} \hat{R}(\vect{\theta}).
\end{equation}
\end{definicion}

Las referencias se citan así~\cite{russell-norvig}.

% ======================= DESARROLLO =======================
\section{Desarrollo}

\subsection{Entorno y herramientas}

\begin{table}[htbp]
  \centering
  \caption{Ejemplo de tabla.}
  \label{tab:entorno}
  \begin{tabular}{ll}
    \toprule
    \textbf{Componente} & \textbf{Versión} \\
    \midrule
    Python & 3.13 \\
    scikit-learn & 1.6 \\
    \bottomrule
  \end{tabular}
\end{table}

\subsection{Implementación}

\begin{lstlisting}[style=python, caption={Ejemplo de código Python.}, label={lst:ejemplo}]
from sklearn.cluster import KMeans

modelo = KMeans(n_clusters=3, random_state=0)
etiquetas = modelo.fit_predict(X)  # asigna cada punto a un grupo
\end{lstlisting}

\begin{lstlisting}[style=bash]
pip install -r requirements.txt
python main.py --datos datos.csv
\end{lstlisting}

\begin{algorithm}[htbp]
\caption{Ejemplo de pseudocódigo}\label{alg:ejemplo}
\begin{algorithmic}[1]
\Require conjunto de datos $D$, número de iteraciones $T$
\Ensure modelo entrenado $\vect{\theta}$
\State inicializar $\vect{\theta}$
\For{$t = 1$ \textbf{hasta} $T$}
  \State $\vect{\theta} \gets \vect{\theta} - \eta \nabla \hat{R}(\vect{\theta})$
\EndFor
\State \Return $\vect{\theta}$
\end{algorithmic}
\end{algorithm}

\begin{aviso}
Algo a tener en cuenta (un error típico, una limitación...).
\end{aviso}

% ======================= RESULTADOS =======================
\section{Resultados}

% Si img/resultado.png no existe aparece un recuadro con el nombre.
\captura{resultado.png}{Ejemplo de captura o gráfica de resultados.}

Como se ve en la Figura~\ref{fig:resultado.png} y la Tabla~\ref{tab:entorno}\ldots

\begin{nota}
Una observación sobre los resultados.
\end{nota}

% ======================= CONCLUSIONES =======================
\section{Conclusiones}

Ideas principales que deja el trabajo, problemas encontrados y posibles
mejoras.

% ======================= USO DE IA =======================
% Se coloca aquí, antes de la bibliografía. Si la quitas, la clase la
% añade igualmente al final del documento.
\declaracionIA

% ======================= REFERENCIAS =======================
\clearpage
\phantomsection
\addcontentsline{toc}{section}{Referencias}
\bibliographystyle{unsrt-es}
\bibliography{bibliografia}

\end{document}
